\section{Conclusion}
We introduced HT-CTS, a calibrated Thompson sampling algorithm for contextual bandits with heavy-tailed rewards, built around an online-calibrated truncated-mean surrogate likelihood. We proved a regret bound of order $\widetilde{O}(d\,T^{1/(1+\epsilon)})$ that interpolates smoothly between the sub-Gaussian rate and the fundamental limits of learning under minimal moment assumptions, and we showed empirically that the resulting algorithm outperforms Gaussian Thompson sampling and optimism-based baselines by a growing margin as tails become heavier. We believe calibrated truncation is a lightweight, broadly applicable modification for deploying randomized-exploration bandit algorithms in environments where reward tails cannot be assumed away.

\bibliographystyle{plainnat}
\bibliography{\jobname}
